\documentclass{article}
\usepackage[utf8]{inputenc}
\usepackage{amsmath}
\usepackage{amssymb}
\usepackage{amsthm}
\usepackage{parskip}

\newtheorem{theorem}{Theorem}
\newtheorem{claim}{Claim}
\newtheorem{remark}{Remark}

\title{The endpoint case using $J+J-3J$ as the main final interval}
\author{Ferran Espu\~na}
\date{\today}

\begin{document}

\maketitle

\begin{theorem}\label{thm:open_two_intervals_endpoint_R6}
Let $A \subset \mathbb T$ be a union of two non-empty disjoint open intervals with
\[
    \mu(A)=\frac15.
\]
If
\[
    A+A-3\cdot A \neq \mathbb T,
\]
then, up to translating $A$ and replacing $A$ by $-A$, we have
\[
    A=\left(0,\frac1{10}\right)\cup\left(\frac12,\frac35\right).
\]
In particular, $2\cdot A$ is contained in an interval of length $\frac15$.
\end{theorem}

\begin{proof}
As before, the condition $A+A-3\cdot A=\mathbb T$ is invariant under translating $A$
and replacing $A$ by $-A$. We may therefore assume
\[
    A=I\cup J, \qquad I=(0,\alpha), \qquad J=(x,x+\beta),
\]
where
\begin{align}
    0<\alpha < x < x+\beta <1, \label{cond:order_R6} \\
    \alpha \geq \beta, \label{cond:alpha_geq_beta_R6} \\
    x-\alpha \leq 1-(x+\beta). \label{cond:shorter_gap_R6}
\end{align}
We have
\[
    \alpha+\beta=\frac15,
\]
and~\eqref{cond:shorter_gap_R6} gives
\begin{equation}\label{ineq:x_leq_3a_2b_R6}
    x\leq \frac{1+\alpha-\beta}{2}=3\alpha+2\beta.
\end{equation}
We will prove that $A+A-3\cdot A=\mathbb T$ unless
\begin{equation}\label{eq:exceptional_case_R6}
    \alpha=\beta=\frac1{10}, \qquad x=\frac12.
\end{equation}

This time we keep $J+J-3\cdot J$ as the main final interval and leave
$I+J-3\cdot J$ out of the main covering argument. Define
\begin{align*}
    R_1=I+I-3\cdot I &= (-3\alpha,2\alpha), \\
    R_2=I+J-3\cdot I &= (x-3\alpha,x+\alpha+\beta), \\
    R_3=J+J-3\cdot I &= (2x-3\alpha,2x+2\beta), \\
    R_4=I+I-3\cdot J &= (-3x-3\beta,-3x+2\alpha), \\
    R_6=J+J-3\cdot J &= (-x-3\beta,-x+2\beta).
\end{align*}
All these intervals are in proper form. Clearly each of them is contained in
$A+A-3\cdot A$.

The first part of the proof is exactly the same as before.

\begin{claim}\label{claim:strict_first_inequalities_R6}
If~\eqref{eq:exceptional_case_R6} does not hold, then
\[
    x<5\alpha
    \qquad\text{and}\qquad
    x<4\alpha+\beta.
\]
\end{claim}

\begin{proof}
By~\eqref{ineq:x_leq_3a_2b_R6} and $\alpha\geq\beta$,
\[
    x\leq 3\alpha+2\beta\leq 5\alpha.
\]
Equality forces $\alpha=\beta$ and $x=3\alpha+2\beta$, hence
$\alpha=\beta=\frac1{10}$ and $x=\frac12$. Thus, outside
\eqref{eq:exceptional_case_R6}, we have $x<5\alpha$.

Similarly,
\[
    x\leq 3\alpha+2\beta\leq 4\alpha+\beta,
\]
and equality again forces~\eqref{eq:exceptional_case_R6}. Hence $x<4\alpha+\beta$.
\end{proof}

\begin{claim}\label{claim:first_cover_R6}
Assume~\eqref{eq:exceptional_case_R6} does not hold. Then
\[
    R_1\cup R_2\cup R_3
\]
covers $[0,2x+2\beta)$ and also covers $(1-3\alpha,1)$.
\end{claim}

\begin{proof}
Modulo $1$, the interval $R_1=(-3\alpha,2\alpha)$ covers
\[
    [0,2\alpha)\cup(1-3\alpha,1).
\]
By Claim~\ref{claim:strict_first_inequalities_R6}, $x<5\alpha$, so $R_2$ contains
$2\alpha$. Therefore $R_1\cup R_2$ covers $[0,x+\alpha+\beta)$.
Again by Claim~\ref{claim:strict_first_inequalities_R6}, $x<4\alpha+\beta$, so
$R_3$ contains $x+\alpha+\beta$. Therefore $R_1\cup R_2\cup R_3$ covers
$[0,2x+2\beta)$.
\end{proof}

Thus the only possible leftover is
\begin{equation}\label{eq:leftover_R6}
    [2x+2\beta,1-3\alpha].
\end{equation}
If this interval is empty, we are done. So assume it is non-empty. Equivalently,
\begin{equation}\label{ineq:two_x_bound_R6}
    2x\leq 2\alpha+3\beta.
\end{equation}
We now cover this leftover using $R_4\cup R_6$.

\begin{claim}\label{claim:R4_R6_cover_except_fake}
The intervals $R_4$ and $R_6$ cover the leftover interval~\eqref{eq:leftover_R6}, except
possibly in the case
\begin{equation}\label{eq:fake_case_R6}
    \alpha=\frac9{65}, \qquad \beta=\frac4{65}, \qquad x=\frac3{13}.
\end{equation}
In that case, the only point of~\eqref{eq:leftover_R6} not covered by $R_4\cup R_6$ is
$\frac{38}{65}$.
\end{claim}

\begin{proof}
Shift the leftover interval by $-1$. We need to cover
\[
    [2x+2\beta-1,-3\alpha].
\]
The left endpoint is strictly to the right of the left endpoint of $R_4$, because
\[
    -3x-3\beta<2x+2\beta-1
\]
is equivalent to $x+\beta>\frac15$, which follows from $x>\alpha$ and
$\alpha+\beta=\frac15$.

The right endpoint is strictly to the left of the right endpoint of $R_6$, because
\[
    -3\alpha<-x+2\beta
\]
is equivalent to $x<3\alpha+2\beta$. This is strict under the present assumption that
\eqref{eq:leftover_R6} is non-empty: indeed, equality in
$x\leq 3\alpha+2\beta$ together with~\eqref{ineq:two_x_bound_R6} would give
$6\alpha+4\beta\leq 2\alpha+3\beta$, impossible.

It remains only to check that $R_4$ and $R_6$ do not leave a gap inside this shifted
leftover interval. Their adjacent endpoints are
\[
    -3x+2\alpha
    \qquad\text{and}\qquad
    -x-3\beta.
\]
By~\eqref{ineq:two_x_bound_R6},
\[
    -x-3\beta\leq -3x+2\alpha.
\]
Thus $R_4$ and $R_6$ overlap, except possibly when
\[
    2x=2\alpha+3\beta.
\]
If this inequality is strict, there is no gap and the whole leftover interval is covered.

Suppose then that $2x=2\alpha+3\beta$. In this case the leftover interval
\eqref{eq:leftover_R6} is a single point, namely $1-3\alpha$. The only point not covered
by $R_4\cup R_6$ is the common endpoint
\[
    -3x+2\alpha=-x-3\beta.
\]
This coincides with $(1-3\alpha)-1=-3\alpha$ precisely when
\[
    -3x+2\alpha=-3\alpha,
\]
or equivalently $3x=5\alpha$. Combining this with $2x=2\alpha+3\beta$ gives
\[
    4\alpha=9\beta.
\]
Since $\alpha+\beta=\frac15$, this gives
\[
    \alpha=\frac9{65}, \qquad \beta=\frac4{65}, \qquad x=\frac3{13}.
\]
In that case the leftover point is
\[
    1-3\alpha=\frac{38}{65}.
\]
This proves the claim.
\end{proof}

The fake case~\eqref{eq:fake_case_R6} is not a genuine solution. It is exactly the case
where the strategy using only $R_4\cup R_6$ leaves an artificial endpoint. The omitted
interval
\[
    R_5=I+J-3\cdot J=(-2x-3\beta,\alpha+\beta-2x)
\]
contains it, because in the fake case
\[
    R_5=\left(-\frac{42}{65},-\frac{17}{65}\right)
\]
and
\[
    \frac{38}{65}-1=-\frac{27}{65}\in R_5.
\]
Thus this point belongs to $A+A-3\cdot A$ after all.

We conclude that, outside the exceptional case~\eqref{eq:exceptional_case_R6}, the set
$A+A-3\cdot A$ is all of $\mathbb T$. Therefore, if
$A+A-3\cdot A\neq\mathbb T$, we must be in the exceptional case
\eqref{eq:exceptional_case_R6}.

In that case,
\[
    A=\left(0,\frac1{10}\right)\cup\left(\frac12,\frac35\right),
\]
and
\[
    2\cdot A=\left(0,\frac15\right),
\]
which is contained in an interval of length $\frac15$. Directly computing the six interval
pieces shows that this exceptional set satisfies
\[
    A+A-3\cdot A=\mathbb T\setminus\left\{\frac15,\frac7{10}\right\}.
\]
This completes the proof.
\end{proof}

\begin{remark}
The proof above shows that one cannot literally omit $I+J-3\cdot J$ from a covering
proof in the open-interval setting. The five intervals $R_1,R_2,R_3,R_4,R_6$ leave a fake
endpoint in the parameter choice
\[
    \alpha=\frac9{65}, \qquad \beta=\frac4{65}, \qquad x=\frac3{13},
\]
and that endpoint is covered exactly by the omitted term $I+J-3\cdot J$. Thus $R_6$ can
be used as the main replacement for $R_5$, but $R_5$ still has to be mentioned to dismiss
this one artificial endpoint.
\end{remark}

\end{document}
